% !TEX root = ../jlt.tex
\section{The Check}
\label{sec:check}

\Cpp{} holds a program to part of the three rules, and a tool holds it to more.
The repository of this paper has such a tool, a linter built on Clang's own libraries, and this section says what it does today.
% TODO(check): this section describes the linter as it is.  To be built, and
% then described here: a strict mode in which every construct is kernel or a
% listed derived form; the modes "check" (all rules satisfied, the axioms
% listed) and "warn" (the axioms and each violation listed) on a named file;
% the three conditions on the order of declarations, conservatively, and across
% files by location; its output on the three listings (concat1 and
% guessed_power refused, the false axiom on the list); a list of the derived
% forms a project licenses beyond the paper's; what it does not see.  None of
% the three listings is under the linter yet: their code is in a probe.

It reads a source file as the compiler does and applies a set of rules to a named part of it.
A rule either refuses the file or reports what it found.
For the first rule it checks, among others, that functions are \cppinline{constexpr}, that call operators are \cppinline{const}, that nothing is thrown and that no mutable state is global.
For the second it checks that a template is constrained and that a generic body calls on its type parameters only what its constraints name, by the names of the calls.
For the third it checks that there is no loop without a stated end and no function that calls itself.

It also prints a list.
Every entry of a trait at one type whose value is \cppinline{true} is an axiom, and the tool lists them with their places, for a file or for a whole library.
The build compares that list with one that is kept under version control, so an axiom cannot be added without the change being seen.

What the tool does not yet do is decide whether a program is Jlt.
That needs a mode in which every construct is either in the kernel or a listed derived form and anything else is refused, and it needs the conditions on the order of declarations (Section~\ref{sec:limits}).
Until then the first clause of this paper's promise, that a template means what its signature says, rests on the programmer and on a reviewer.
